% ============================================================================
% REGENERATED, 2026-09-19, against v2 (irregular-series, chronological-
% primary-split) numbers -- Gate v2 (M2_Outcome_Definition_Gate.md),
% signed off 2026-09-19 by Akram. Replaces the prior version, which
% predated Deviation 7 (2026-08-29 fixed-T=2 cohort) and was STALE-
% bannered throughout. Source of truth for every number here:
% M2_Findings_Dossier.md Sec 24 (v2, primary) and Secs 1-22 (T=2,
% secondary/historical, retained as \S\ref{sec:t2-legacy}).
%
% SKELETON ONLY, per standing project policy: real section structure,
% real tables (\input from paper_assets/tables/*.tex, each independently
% regenerable from a saved data file -- see each table's own header
% comment), real figures, real citations -- but NO argumentative prose.
% Every narrative/discussion subsection is an explicit placeholder
% pointing at the dossier section it will be written from, deferred
% until all of M1-M5 reach Green-core status and manuscript writing
% formally begins across the whole SihhaDZ portfolio.
% ============================================================================
\documentclass[11pt]{article}

% ---------- Packages ----------
\usepackage[a4paper,margin=2.5cm]{geometry}
\usepackage[utf8]{inputenc}
\usepackage[T1]{fontenc}
\usepackage{lmodern}
\usepackage{booktabs}
\usepackage{graphicx}
\usepackage{amsmath,amssymb}
\usepackage{xcolor}
\usepackage{caption}
\usepackage{subcaption}
\usepackage{array}
\usepackage{longtable}
\usepackage{float}
\usepackage{enumitem}
\usepackage[most]{tcolorbox}
\usepackage{fancyhdr}
\usepackage[colorlinks=true,linkcolor=accentdark,citecolor=accentdark,urlcolor=accentdark]{hyperref}

% ---------- Colors (matching the project's established report palette) ----------
\definecolor{ink}{HTML}{1B2430}
\definecolor{inksoft}{HTML}{4E5964}
\definecolor{accent}{HTML}{1C6E68}
\definecolor{accentdark}{HTML}{123F3B}
\definecolor{accentsoft}{HTML}{E1EEEC}
\definecolor{amber}{HTML}{93591F}
\definecolor{ambersoft}{HTML}{F3E7D9}
\definecolor{brick}{HTML}{8E3333}
\definecolor{bricksoft}{HTML}{F5E2E1}

\newtcolorbox{finding}[1]{colback=bricksoft,colframe=brick,fonttitle=\bfseries,title=#1,left=2mm,right=2mm,top=1mm,bottom=1mm}
\newtcolorbox{honest}[1]{colback=ambersoft,colframe=amber,fonttitle=\bfseries,title=#1,left=2mm,right=2mm,top=1mm,bottom=1mm}
\newtcolorbox{note}[1]{colback=accentsoft,colframe=accent,fonttitle=\bfseries,title=#1,left=2mm,right=2mm,top=1mm,bottom=1mm}
\newtcolorbox{prosetodo}[1]{colback=gray!8,colframe=gray!50,fonttitle=\bfseries,title=#1,left=2mm,right=2mm,top=1mm,bottom=1mm}

\pagestyle{fancy}
\fancyhf{}
\lhead{\small\textcolor{inksoft}{SihhaDZ --- Module 2}}
\rhead{\small\textcolor{inksoft}{Clinical Risk Prediction (v2, Gate-signed-off)}}
\cfoot{\small\thepage}
\renewcommand{\headrulewidth}{0.4pt}

\newcommand{\good}[1]{\textcolor{accentdark}{\textbf{#1}}}
\newcommand{\bad}[1]{\textcolor{brick}{\textbf{#1}}}

\title{\textbf{\Large Module 2 --- Clinical Risk Prediction:\\[2pt] GRU-D, Adaptive Conformal Inference, and the\\[2pt] Irregular-Series Hybrid Temporal--Static Controlled Experiment}\\[8pt]
\large Manuscript Skeleton (Structure and Results Only --- Prose Deferred)}
\author{Student 2 --- M2, SihhaDZ Emergency Platform Research Program}
\date{Regenerated 19 September 2026 (Gate v2 results)}

\begin{document}
\maketitle

\begin{center}
\begin{tabular}{ll}
\textbf{Dataset:} & Full credentialed MIMIC-IV-ED 2.2 (418{,}025 gate-eligible ED visit rows) \\
\textbf{Track:} & Track A (Research) --- SihhaDZ Dual-Track Contract \\
\textbf{Target venue:} & IEEE JBHI (SihhaDZ roadmap primary for M2) \\
\textbf{Primary split:} & Chronological (rolling-origin), cut at 2016 --- Gate v2 \S7 \\
\textbf{Companion assets:} & \texttt{M2\_Findings\_Dossier.md} \S24, \texttt{M2\_Hybrid\_Deviation\_Log.md}, \texttt{M2\_Outcome\_Definition\_Gate.md} (v2) \\
\end{tabular}
\end{center}

\vspace{4mm}
\begin{note}{Golden Rule}
Every number in this report is read from a saved result file produced by an executed script. No metric is estimated, extrapolated, or invented.
\end{note}

\begin{prosetodo}{Prose deferred}
This document is a structural skeleton: real section headings, real tables (each \texttt{\textbackslash input}-ed from an independently regenerable \texttt{paper\_assets/tables/*.tex} file), real figures, and real, verified citations -- but no argumentative or narrative prose. Per standing project policy, manuscript prose is written only once all of M1--M5 reach Green-core status, so that framing decisions are made once, consistently, across the whole portfolio, not module-by-module. Each placeholder below names the \texttt{M2\_Findings\_Dossier.md} section its eventual prose will draw from.
\end{prosetodo}

\tableofcontents
\newpage

% ======================================================================
\section{Executive Summary}
% ======================================================================
\begin{prosetodo}{Deferred}
Draws from \texttt{M2\_Findings\_Dossier.md} \S24 (v2 results) and \S23 (external Q1-readiness review, four blockers, all fixed/closed 2026-09-19).
\end{prosetodo}

% ======================================================================
\section{Introduction and Research Questions}
% ======================================================================

Per \texttt{M2\_Outcome\_Definition\_Gate.md} (v2, signed off 2026-09-19), this module answers three questions under a single, identically-defined cohort, outcome, and split across every comparator:

\begin{enumerate}[leftmargin=*]
\item \emph{Mandatory core:} given a stream of ED vital-sign observations, can GRU-D (\cite{che2018recurrent}) with MAPIE-calibrated prediction intervals outperform a clinical early-warning score (NEWS2)?
\item \emph{Primary research question (pre-registered comparison):} under real temporal non-stationarity (the chronological primary split), are online conformal updates (Error-quantified Conformal Inference, ECI, \cite{wu2025error}; and an online class-conditional variant) needed to keep marginal and class-conditional (Mondrian) coverage near nominal for the clinically-relevant minority (deterioration) class, or does periodic recalibration of the corresponding static split-conformal sets suffice?
\item \emph{Controlled experiment (supervisor-commissioned):} under strictly matched clinical information, does representing vital-sign data as an irregular temporal sequence (GRU-D) add predictive value over a non-temporal, tabular aggregation (XGBoost) -- and is any such value attributable to the temporal representation itself, or to the static-feature fusion architecture around it?
\end{enumerate}

\begin{prosetodo}{Deferred}
Full motivation and related-work framing draws from \texttt{M2\_Findings\_Dossier.md} \S1 and \texttt{M2\_Related\_Work\_Notes.md}.
\end{prosetodo}

% ======================================================================
\section{Data, Cohort, and Outcome Definition}
\label{sec:data}
% ======================================================================

\subsection{Cohort}
Full credentialed MIMIC-IV-ED 2.2, filtered to age~$\geq$~18 and non-null triage acuity: \textbf{418{,}025 rows}. The outcome, \texttt{outcome\_critical}, is a composite critical-outcome label (in-hospital mortality \textbf{OR} direct ED-to-ICU transfer within 12h of ED disposition), locked in Gate v2 \S2 identically across every model in this report.

\subsection{Primary split: chronological (rolling-origin)}
Gate v2 \S7 promotes the chronological split to primary (cut year 2016), reversing the T=2 pipeline's random-split convention, to genuinely exercise ECI's premise (adaptive calibration under real, not merely simulated, temporal non-stationarity):

\begin{table}[H]
\centering
\caption{Chronological-primary split (Gate v2 \S7), \texttt{m2\_build\_irregular\_cohort.py}}
\begin{tabular}{lrr}
\toprule
Period & Rows & Unique patients \\
\midrule
Train (pre-2016) & 288{,}208 & 139{,}371 \\
Test (post-2016) & 84{,}241 & 61{,}843 \\
Dropped (straddles the 2016 cut) & 45{,}576 & 18{,}968 \\
\bottomrule
\end{tabular}
\end{table}

A secondary random 80/20 split is persisted (\texttt{data/m2\_v2\_random\_split\_ids.json}) as a planned robustness arm, deliberately deferred this session (\texttt{M2\_Findings\_Dossier.md} \S24.5) pending a parametrized-caching refactor that does not risk the primary split's already-validated artifacts.

\subsection{Irregular vital-sign sequence}
Every prior M2 result (Secs~\ref{sec:t2-legacy}) used a fixed two-timestep representation (triage, latest-ED). This report's primary results use the \textbf{full irregular \texttt{vitalsign.csv} sequence}: triage plus up to $T_{\max}=10$ most-recent ED observations per stay ($T_{\max}$ = the measured 90th percentile of ED-observation count; 0 stays truncated in either split period). A genuine, previously-undocumented recording-practice shift was found and confound-checked (\texttt{M2\_Findings\_Dossier.md} \S24.1/\S24.5): 4.62\% of train-period (pre-2016) stays have zero recorded ED vitals vs.\ 0.02\% of test-period (post-2016) stays; stratifying every model's AUROC by ED-observation presence found no evidence this confounds any discrimination result.

\subsection{Prediction horizon}
Gate v2 \S3 retracts the T=2 pipeline's horizon construction (an ED-length-of-stay filter that was found to select an easier cohort, not a genuine horizon effect -- Sec~\ref{sec:t2-legacy}) and replaces it with a real time-to-event label $y_h = \mathbb{1}\{0 \le \text{event\_time} - \text{origin} \le h\}$, $h \in \{1\text{h}, 6\text{h}, 12\text{h}\}$, measured on the full eligible cohort ($n=417{,}964$): positive rates 0.055\% (1h, 228 positives), 3.391\% (6h, 14{,}175), 5.152\% (12h, 21{,}534). Every result in Sections~\ref{sec:results}.1--\ref{sec:results}.4 uses the un-horizoned \texttt{outcome\_critical} composite as the training label; the horizons 1h/6h/12h ablation CLAUDE.md's M2 directive requires is reported separately in Sec.~\ref{sec:horizon} on these $y_h$ labels.

% ======================================================================
\section{Methods}
% ======================================================================

\subsection{Six compared models}
\begin{table}[H]
\centering
\caption{The six v2 models -- matched information, different representations}
\begin{tabular}{p{3.0cm}p{3.4cm}p{2.6cm}p{5.0cm}}
\toprule
Model & Temporal input & Static input & Role \\
\midrule
Static-Only MLP & --- & 4 statics & No-vitals floor \\
mTAN (\cite{shukla2021multi}) & full irregular sequence & --- & Modern attention-based comparator (Q1-blocker \#1 fix) \\
Temporal-Only GRU-D & full irregular sequence & --- & Clean temporal baseline \\
Frozen-Temporal Head Control & frozen encoder (above) & --- & Isolates static branch's contribution \\
XGBoost-Matched & aggregated (mean/min/max/last/count) & 4 statics & Matched-information tabular representation \\
\textbf{Hybrid GRU-D} & frozen encoder & 4 statics & \textbf{Primary model} \\
\bottomrule
\end{tabular}
\end{table}

\subsection{Tuning symmetry}
Every model in the table above received a real, comparable search budget before its 5-seed confirmatory run: GRU-D-family/XGBoost via a 3-seed, $\sim$14--15-candidate staged search (anti-p-hacking \texttt{UNTUNED\_BASE} re-merge, baseline-vs-SD noise guard, ECE calibration guard); mTAN, after an initial lighter 1-seed protocol, was re-tuned at matched rigor on 2026-09-20 (3 seeds, 4 stages, same guards; \texttt{M2\_Hybrid\_Deviation\_Log.md} Deviation 15, resolved). XGBoost-Matched v2 and the two head-only architectures (Frozen-Temporal Head Control, Hybrid GRU-D) were tuning-symmetry-resolved on 2026-09-19 (Deviations 14, 16), replacing earlier disclosed reuse/untuned configs.

\subsection{Conformal methods}
MAPIE (split-conformal, \cite{mapie}), ECI (\cite{wu2025error}), and Mondrian-LAC (label-conditional split-conformal sets, \cite{sadinle2019least}) are compared at $\alpha=0.1$ (90\% nominal coverage) on the three GRU-D-family models, reporting marginal and class-1 (deterioration) coverage separately (Q1-blocker \#2 fix).

\subsection{Inferential statistics}
Every comparison uses a \textbf{patient-cluster bootstrap} (10{,}000 resamples/seed, patients resampled with replacement, all of a sampled patient's encounters retained with multiplicity, never deduplicated), with Holm step-down correction across the 5-delta confirmatory family. DeLong's test is omitted per the project's sanctioned fallback (no validated patient-cluster-adjusted implementation exists).

% ======================================================================
\section{Results}
\label{sec:results}
% ======================================================================

\subsection{Headline discrimination}

\begin{figure}[H]
\centering
\includegraphics[width=0.85\textwidth]{figures/fig1_headline_auroc_v2.pdf}
\caption{All 6 v2 models, chronological-primary split, mean $\pm$ SD across 5 seeds.}
\end{figure}

\begin{table*}[t]
\centering
\caption{Discrimination and calibration, all 6 models, primary cohort, chronological split. Three models (XGBoost-Matched, Frozen-Temporal Head Control, Hybrid GRU-D) are reported at their final tuned configurations. Mean $\pm$ SD over five training seeds. $^{\ddagger}$XGBoost-Matched is deterministic (identical predictions for every seed), so its cells give the point value and the 95\% patient-cluster bootstrap interval (10,000 resamples), a measure of population uncertainty rather than seed variability. Best value per column in bold (highest AUROC and AUPRC, lowest Brier score and ECE). Proxy cohort (MIMIC-IV-ED-derived).}
\label{tab:m2-v2-headline}
\resizebox{\textwidth}{!}{%
\begin{tabular}{lcccc}
\toprule
Model & AUROC & AUPRC & Brier & ECE \\
\midrule
Static-Only MLP & 0.8556 $\pm$ 0.0007 & 0.3596 $\pm$ 0.0022 & 0.0638 & 0.0334 \\
mTAN (tuned) & 0.8622 $\pm$ 0.0031 & 0.4996 $\pm$ 0.0082 & 0.0549 & 0.0203 \\
Temporal-Only GRU-D & 0.9209 $\pm$ 0.0022 & 0.6638 $\pm$ 0.0039 & 0.0431 & \textbf{0.0077} \\
Frozen-Temporal Head Control (tuned) & 0.9201 $\pm$ 0.0012 & 0.6650 $\pm$ 0.0035 & 0.0435 & 0.0153 \\
XGBoost-Matched (tuned) & 0.9271 [0.9241, 0.9301]$^{\ddagger}$ & 0.6337 [0.6221, 0.6454]$^{\ddagger}$ & 0.0464 & 0.0176 \\
Hybrid GRU-D (primary, tuned) & \textbf{0.9372 $\pm$ 0.0006} & \textbf{0.6873 $\pm$ 0.0015} & \textbf{0.0418} & 0.0139 \\
\bottomrule
\end{tabular}%
}
\end{table*}


\subsection{Calibration}

\begin{figure}[H]
\centering
\includegraphics[width=0.85\textwidth]{figures/fig2_calibration_v2.pdf}
\caption{Reliability diagrams, all 6 v2 models.}
\end{figure}

\subsection{The five primary comparisons: patient-cluster bootstrap}
\label{sec:deltas}

\begin{figure}[H]
\centering
\includegraphics[width=0.85\textwidth]{figures/fig4_deltas_forest_v2.pdf}
\caption{The five v2 deltas, patient-cluster bootstrap 95\% CIs, Holm-corrected. All five are significant at $\alpha=0.05$.}
\end{figure}

\begin{table*}[t]
\centering
\caption{The five pre-specified pairwise differences in AUROC (first model minus second), patient-cluster bootstrap 95\% CIs, Holm-corrected for multiplicity. All values use the final tuned configurations. Every difference is significant at $\alpha=0.05$. The difference between Temporal-Only GRU-D and XGBoost-Matched does not, on its own, isolate temporal representation, because only XGBoost-Matched receives the static features (Sec.~\ref{sec:results-discrimination}); on the composite outcome it shows tuned XGBoost-Matched significantly ahead of Temporal-Only GRU-D alone. Significant differences in bold. Proxy cohort (MIMIC-IV-ED-derived).}
\label{tab:m2-v2-deltas}
\resizebox{\textwidth}{!}{%
\begin{tabular}{lcccc}
\toprule
Comparison & Mean & SD (seeds) & Holm-corrected $p$ & Significant \\
\midrule
Hybrid GRU-D $-$ XGBoost-Matched (primary comparison) & \textbf{+0.0101} & 0.0006 & $<$0.001 & Yes \\
Hybrid GRU-D $-$ Frozen-Temporal Head Control (static fusion) & \textbf{+0.0171} & 0.0007 & $<$0.001 & Yes \\
Temporal-Only GRU-D $-$ XGBoost-Matched & \textbf{-0.0063} & 0.0022 & 0.0022 & Yes \\
Static-Only MLP $-$ Temporal-Only GRU-D & \textbf{-0.0653} & 0.0025 & $<$0.001 & Yes \\
Hybrid GRU-D $-$ mTAN & \textbf{+0.0750} & 0.0034 & $<$0.001 & Yes \\
\bottomrule
\end{tabular}%
}
\end{table*}


\begin{prosetodo}{Deferred}
Draws from \texttt{M2\_Findings\_Dossier.md} \S24.2b, including the $\Delta_{\text{temporal\_alone}}$ sign-reversal finding under tuning symmetry.
\end{prosetodo}

\subsection{Conformal coverage: MAPIE vs.\ ECI vs.\ Mondrian-LAC}
\label{sec:conformal}

\begin{figure}[H]
\centering
\includegraphics[width=0.85\textwidth]{figures/fig3_conformal_coverage_v2.pdf}
\caption{Marginal and class-1 coverage, MAPIE vs.\ ECI vs.\ Mondrian-LAC, chronological-primary split -- this module's central result (\texttt{M2\_Findings\_Dossier.md} \S24.3).}
\end{figure}

\begin{table*}[t]
\centering
\caption{Marginal and class-1 (deterioration) coverage of MAPIE (split-conformal), ECI (adaptive) and Mondrian-LAC (label-conditional) prediction sets on the chronological primary split; the nominal level is 90\%. Entries are five-seed means, with the standard deviation shown for Mondrian-LAC. MAPIE here is the primary classification-set run (MAPIE library); the static split-conformal column of Table~\ref{tab:m2-eci-vs-recal} reports a separately recomputed interval construction of the same quantile (see its caption). Class-1 coverage values at or above the validity threshold of 0.89 are in bold. Proxy cohort (MIMIC-IV-ED-derived).}
\label{tab:m2-v2-conformal}
\resizebox{\textwidth}{!}{%
\begin{tabular}{lccc}
\toprule
Model & MAPIE (marginal / class-1) & ECI (marginal / class-1) & Mondrian-LAC (class-1) \\
\midrule
Temporal-Only GRU-D & 0.8984 / 0.1903 & 0.9003 / 0.2153 & \textbf{0.8939 $\pm$ 0.0069} \\
Frozen-Temporal Head Control & 0.9003 / 0.1860 & 0.8996 / 0.2104 & 0.8885 $\pm$ 0.0064 \\
Hybrid GRU-D (primary) & 0.9004 / 0.2043 & 0.8998 / 0.2227 & \textbf{0.8951 $\pm$ 0.0050} \\
\bottomrule
\end{tabular}%
}
\end{table*}


\begin{table*}[t]
\centering
\caption{Clinician-facing view of the coverage results in Table~\ref{tab:m2-v2-conformal}. Among patients who deteriorate (class 1): how often each method's 90\%-target prediction set or interval \emph{fails} to cover the outcome (miss rate, the falsely reassuring case), how often it abstains (empty-set rate), and the average interval width or set size. MAPIE's width and set-size cell is left blank (--) because its interval width was not persisted per class. Entries are five-seed means $\pm$ SD. For each model, the method with the highest class-1 coverage (lowest miss rate) is in bold. Proxy cohort (MIMIC-IV-ED-derived).}
\label{tab:m2-g7-clinician-decision}
\resizebox{\textwidth}{!}{%
\begin{tabular}{llcccc}
\toprule
Model & Method & Coverage (class 1, \%) & Miss rate (class 1, \%) & Empty-set rate (class 1, \%) & Width / set size \\
\midrule
Temporal-Only GRU-D & MAPIE (split-conformal) & 19.0 $\pm$ 2.0 & 81.0 $\pm$ 2.0 & 45.6 $\pm$ 1.4 & -- \\
 & ECI (adaptive) & 21.5 $\pm$ 2.6 & 78.5 $\pm$ 2.6 & -- & 0.316 $\pm$ 0.013 \\
 & Mondrian-LAC (label-conditional) & \textbf{89.4 $\pm$ 0.7} & \textbf{10.6 $\pm$ 0.7} & 0.0 $\pm$ 0.0 & 1.128 $\pm$ 0.017 \\
Frozen-Temporal Head Control & MAPIE (split-conformal) & 18.6 $\pm$ 0.7 & 81.4 $\pm$ 0.7 & 45.5 $\pm$ 0.9 & -- \\
 & ECI (adaptive) & 21.0 $\pm$ 0.8 & 79.0 $\pm$ 0.8 & -- & 0.268 $\pm$ 0.010 \\
 & Mondrian-LAC (label-conditional) & \textbf{88.8 $\pm$ 0.6} & \textbf{11.2 $\pm$ 0.6} & 0.0 $\pm$ 0.0 & 1.125 $\pm$ 0.013 \\
Hybrid GRU-D (primary) & MAPIE (split-conformal) & 20.4 $\pm$ 0.9 & 79.6 $\pm$ 0.9 & 46.0 $\pm$ 0.8 & -- \\
 & ECI (adaptive) & 22.3 $\pm$ 1.1 & 77.7 $\pm$ 1.1 & -- & 0.258 $\pm$ 0.010 \\
 & Mondrian-LAC (label-conditional) & \textbf{89.5 $\pm$ 0.5} & \textbf{10.5 $\pm$ 0.5} & 0.0 $\pm$ 0.0 & 1.091 $\pm$ 0.008 \\
\bottomrule
\end{tabular}%
}
\end{table*}


\begin{prosetodo}{Deferred}
Draws from \texttt{M2\_Findings\_Dossier.md} \S24.3/\S24.4, including the precision note that ECI's near-nominal coverage is marginal only (class-1 coverage 0.05--0.08) and only Mondrian-LAC achieves class-conditional validity -- the two corrections are orthogonal and both required.
\end{prosetodo}

\subsection{Horizon ablation}
\label{sec:horizon}
\begin{table*}[t]
\centering
\caption{Sensitivity to the outcome definition: the same models and the same full per-stay input sequence, evaluated against three time-anchored label definitions $y_h$, $h\in\{1\mathrm{h}, 6\mathrm{h}, 12\mathrm{h}\}$, on the chronological split. The input is not truncated at $h$, so the rows compare outcome definitions, not predictions made $h$ hours ahead. Encoders were trained per definition with the hyperparameters tuned for the composite outcome; 3 of 6 models. ``Sig.'' requires Holm significance and all 5 per-seed confidence intervals excluding zero. $^{\dagger}$Underpowered (few test and validation positives); not interpreted. Highest AUROC per definition and significant differences in bold. Proxy cohort (MIMIC-IV-ED-derived).}
\label{tab:m2-horizon}
\resizebox{\textwidth}{!}{%
\begin{tabular}{lrccccc}
\toprule
Label definition & Test pos. & Hybrid & Temporal-Only & XGBoost & Hybrid $-$ XGB (95\% CI; sig.) & Temporal $-$ XGB (95\% CI; sig.) \\
\midrule
1h$^{\dagger}$ & 66 & \textbf{0.9311 $\pm$ 0.0051} & 0.8761 $\pm$ 0.0585 & 0.9205 $\pm$ 0.0000 & +0.0106 [-0.0088, +0.0304]; no & -0.0445 [-0.0885, -0.0034]; no \\
6h & 3,841 & \textbf{0.9704 $\pm$ 0.0010} & 0.9624 $\pm$ 0.0016 & 0.9463 $\pm$ 0.0000 & \textbf{+0.0241 [+0.0221, +0.0262]; yes} & \textbf{+0.0161 [+0.0135, +0.0187]; yes} \\
12h & 5,875 & \textbf{0.9506 $\pm$ 0.0004} & 0.9370 $\pm$ 0.0019 & 0.9356 $\pm$ 0.0000 & \textbf{+0.0151 [+0.0134, +0.0168]; yes} & +0.0014 [-0.0012, +0.0040]; no \\
\bottomrule
\end{tabular}%
}
\end{table*}


\begin{prosetodo}{Deferred}
Draws from \texttt{M2\_Findings\_Dossier.md} \S24.6 and Deviation 20. Key point for the eventual prose: the sign of the temporal-only effect is outcome-dependent, so \S24.2b's composite-outcome finding must not be stated as general.
\end{prosetodo}

\subsection{Pre-registered comparison against cheap remedies}
\label{sec:cheap-remedies}
\begin{table*}[t]
\centering
\caption{Static split-conformal, ECI and periodic recalibration of split-conformal: marginal coverage and mean clipped interval width (mean $\pm$ SD over seeds) on the chronological test period. Valid means marginal coverage $\geq 0.89$. ECI is run as implemented in the authors' reference code. Labels are assumed immediate for all methods. The static column is an independently recomputed interval $[\hat{p}-q,\hat{p}+q]\cap[0,1]$ built from the same LAC quantile as the MAPIE classification-set run of Table~\ref{tab:m2-v2-conformal}, so that its width is comparable with the other three columns. The pipeline requires the two to agree in marginal coverage to within 0.003 on every model and seed (largest gap between the unrounded means: 0.0001). For each model, the valid method with the narrowest interval is in bold. Proxy cohort (MIMIC-IV-ED-derived).}
\label{tab:m2-eci-vs-recal}
\resizebox{\textwidth}{!}{%
\begin{tabular}{lcccccccc}
\toprule
Model & \multicolumn{2}{c}{static split-conformal} & \multicolumn{2}{c}{ECI} & \multicolumn{2}{c}{recal.\ ($W{=}10$k)} & \multicolumn{2}{c}{recal.\ ($W{=}40$k)} \\
\cmidrule(lr){2-3}\cmidrule(lr){4-5}\cmidrule(lr){6-7}\cmidrule(lr){8-9}
 & Cov. & Width & Cov. & Width & Cov. & Width & Cov. & Width \\
\midrule
Hybrid GRU-D (primary) & \textbf{0.9004 $\pm$ 0.0009} & \textbf{0.235 $\pm$ 0.010} & 0.8998 $\pm$ 0.0011 & 0.258 $\pm$ 0.010 & 0.9005 $\pm$ 0.0002 & 0.236 $\pm$ 0.010 & 0.9012 $\pm$ 0.0002 & 0.238 $\pm$ 0.010 \\
Temporal-Only GRU-D & \textbf{0.8984 $\pm$ 0.0006} & \textbf{0.291 $\pm$ 0.019} & 0.9003 $\pm$ 0.0011 & 0.316 $\pm$ 0.013 & 0.9005 $\pm$ 0.0002 & 0.300 $\pm$ 0.020 & 0.9010 $\pm$ 0.0003 & 0.302 $\pm$ 0.019 \\
Frozen-Temporal Head Control & \textbf{0.9004 $\pm$ 0.0003} & \textbf{0.229 $\pm$ 0.014} & 0.8996 $\pm$ 0.0005 & 0.268 $\pm$ 0.010 & 0.9006 $\pm$ 0.0002 & 0.231 $\pm$ 0.013 & 0.9012 $\pm$ 0.0002 & 0.233 $\pm$ 0.013 \\
\bottomrule
\end{tabular}%
}
\end{table*}


\begin{table*}[t]
\centering
\caption{Pre-specified comparators for the online label-conditional construction (B), Hybrid GRU-D (primary model). A method is \emph{valid} if coverage of both classes is $\geq 0.89$. A: static Mondrian-LAC; B: Mondrian-ECI; C1 and C2: static Mondrian with periodic recalibration over the two window lengths shown; D and E: conditioning on bins of the model's own risk score instead of the true label, static (D) and online (E). Labels are assumed immediate for all methods; recalibration windows are in arrivals, not calendar time. The valid method with the smallest mean set size is in bold. Proxy cohort (MIMIC-IV-ED-derived).}
\label{tab:m2-mondrian-comparators}
\resizebox{\textwidth}{!}{%
\begin{tabular}{lccccc}
\toprule
Method & Class-0 cov. & Class-1 cov. & Mean $|C|$ & Informative $P(C{=}\{1\}\mid y{=}1)$ & Valid \\
\midrule
A: static Mondrian-LAC & 0.914 $\pm$ 0.001 & 0.895 $\pm$ 0.005 & 1.091 $\pm$ 0.008 & 0.794 $\pm$ 0.005 & yes \\
B: Mondrian-ECI (online, label-cond.) & 0.900 $\pm$ 0.001 & 0.898 $\pm$ 0.001 & 1.133 $\pm$ 0.039 & 0.808 $\pm$ 0.004 & yes \\
C1: static Mondrian + recal.\ ($W{=}10$k) & 0.901 $\pm$ 0.000 & 0.901 $\pm$ 0.001 & 1.087 $\pm$ 0.003 & 0.816 $\pm$ 0.003 & yes \\
C2: static Mondrian + recal.\ ($W{=}40$k) & \textbf{0.903 $\pm$ 0.000} & \textbf{0.899 $\pm$ 0.001} & \textbf{1.085 $\pm$ 0.004} & \textbf{0.813 $\pm$ 0.004} & yes \\
D: static score-bin Mondrian & 0.927 $\pm$ 0.003 & 0.636 $\pm$ 0.009 & 0.996 $\pm$ 0.005 & 0.266 $\pm$ 0.020 & no \\
E: online score-bin ECI & 0.956 $\pm$ 0.006 & 0.648 $\pm$ 0.008 & 1.030 $\pm$ 0.006 & 0.266 $\pm$ 0.011 & no \\
\bottomrule
\end{tabular}%
}
\end{table*}


\begin{prosetodo}{Deferred}
Draws from \texttt{M2\_ECI\_vs\_Recalibration\_PreRegistration.md}, \texttt{M2\_MondrianECI\_PreRegistration.md} (Parts 1--4), and \texttt{M2\_Findings\_Dossier.md} \S24.3b, \S24.3c and the revision block at the top of \S24.4. Both comparisons were pre-registered before the comparators were run, and both returned NOT NEEDED for the online method.
\end{prosetodo}

\subsection{Label lag: does the verdict survive late outcomes?}
\label{sec:label-lag}
\begin{table*}[t]
\centering
\caption{Label lag (pre-specified), Hybrid GRU-D (primary model). Proxy cohort (MIMIC-IV-ED-derived). Lag $L$ is in arrivals; hour equivalents assume $\sim$5 arrivals/hour and are approximate ($\pm$50\%). \emph{Marginal:} coverage / mean clipped interval width for ECI and for static split-conformal with periodic recalibration ($W$=10k and $W$=40k). Valid means coverage $\geq$ 0.89; ECI is not needed under these conditions if a valid recalibration is at most 0.02 wider. For each lag, the valid method with the narrowest interval is in bold. $L=\infty$ is the setting in which no label is ever observed, so recalibration reduces to the static interval. Realistic band (pre-specified): $L\in\{60, 240, 1000\}$. Labels are otherwise assumed uniform-lag; ECI's learning rate is re-selected per lag on the calibration-selection fold.}
\label{tab:m2-label-lag-marginal}
\resizebox{\textwidth}{!}{%
\begin{tabular}{llccc}
\toprule
$L$ (arrivals) & $\approx$ time & ECI & recal.\ $W$=10k & recal.\ $W$=40k \\
\midrule
0 & 0 & 0.8998 / 0.258 & \textbf{0.9005 / 0.236} & 0.9012 / 0.238 \\
60 & $\approx$12\,h & 0.9020 / 0.245 & \textbf{0.9005 / 0.236} & 0.9012 / 0.238 \\
240 & $\approx$2\,d & 0.9020 / 0.245 & \textbf{0.9005 / 0.236} & 0.9012 / 0.238 \\
1{,}000 & $\approx$8\,d & 0.9021 / 0.250 & \textbf{0.9004 / 0.236} & 0.9012 / 0.238 \\
4{,}000 & $\approx$33\,d & 0.9923 / 0.949 & \textbf{0.9007 / 0.236} & 0.9013 / 0.238 \\
16{,}000 & $\approx$4.4\,mo & 0.9547 / 0.685 & \textbf{0.9009 / 0.237} & 0.9013 / 0.238 \\
$\infty$ & never & 0.9157 / 0.321 & \textbf{0.9004 / 0.235} & \textbf{0.9004 / 0.235} \\
\bottomrule
\end{tabular}%
}
\end{table*}

\begin{table*}[t]
\centering
\caption{Label lag, class-conditional (Mondrian) sets, Hybrid GRU-D (primary model). Proxy cohort (MIMIC-IV-ED-derived). Class-0 / class-1 coverage, mean $|C|$ and the informative rate $P(C{=}\{1\}\mid y{=}1)$ for Mondrian-ECI (B) and static Mondrian with periodic recalibration (C1: $W$=10k). Valid means both coverages $\geq$ 0.89. The static anchor (A) uses no label feedback and so does not depend on lag; its class-0 / class-1 coverage and mean $|C|$ are 0.914 / 0.895 / 1.091 (valid). For each lag, the valid method with the smaller mean set size is in bold.}
\label{tab:m2-label-lag-class}
\resizebox{\textwidth}{!}{%
\begin{tabular}{llcccc}
\toprule
$L$ (arrivals) & $\approx$ time & B: cov$_0$ / cov$_1$ & B: $|C|$ & B: informative & C1: cov$_0$ / cov$_1$ / $|C|$ / informative \\
\midrule
0 & 0 & 0.900 / 0.898 & 1.133 & 0.808 & \textbf{0.901 / 0.901 / 1.087 / 0.816} \\
60 & $\approx$12\,h & 0.901 / 0.899 & 1.095 & 0.813 & \textbf{0.901 / 0.901 / 1.087 / 0.816} \\
240 & $\approx$2\,d & 0.901 / 0.900 & 1.098 & 0.813 & \textbf{0.901 / 0.902 / 1.088 / 0.816} \\
1{,}000 & $\approx$8\,d & 0.901 / 0.900 & 1.199 & 0.793 & \textbf{0.901 / 0.902 / 1.088 / 0.816} \\
4{,}000 & $\approx$33\,d & 0.880 / 0.893 & 1.451 & 0.301 & \textbf{0.901 / 0.902 / 1.088 / 0.815} \\
16{,}000 & $\approx$4.4\,mo & 0.797 / 0.907 & 1.388 & 0.353 & \textbf{0.903 / 0.900 / 1.089 / 0.812} \\
$\infty$ & never & 0.913 / 0.955 & 1.335 & 0.797 & \textbf{0.914 / 0.895 / 1.091 / 0.794} \\
\bottomrule
\end{tabular}%
}
\end{table*}


\begin{prosetodo}{Deferred}
Draws from \texttt{M2\_LabelLag\_PreRegistration.md} and \texttt{M2\_Findings\_Dossier.md} \S24.3d. Pre-registered decision band $L\in\{60,240,1000\}$ arrivals; overall reading recorded in the dossier and the pre-registration's RESULT section. Limits to carry into the prose: uniform lag only, approximate arrival-to-hour conversion, one cohort, two recalibration windows, and label modification by intervention (see \texttt{M2\_Related\_Work\_Notes.md} \S7).
\end{prosetodo}

\subsection{Small required-ablation re-runs (logit vs.\ clipping; streaming latency)}
\label{sec:small-ablations}
\begin{table*}[t]
\centering
\caption{Interval bounding and streaming latency, chronological split (Proxy cohort (MIMIC-IV-ED-derived)). \emph{Logit vs.\ clipping:} marginal coverage / class-1 coverage / mean width of the same ECI radius under each bounding, 5-seed means. The logit-bounded interval lies strictly inside $(0,1)$ and so cannot contain a hard label: its coverage is 0 by construction. \emph{Latency:} batch size 1, 2,000 timed calls, Hybrid GRU-D (primary model) seed 42, NVIDIA GeForce RTX 3050 and 8 logical CPUs; hardware- and load-dependent, this machine only; the production deployment path is not measured here. Budget: 2,000\,ms per patient update. For each model, the bounding with the higher coverage is in bold.}
\label{tab:m2-small-ablations}
\resizebox{\textwidth}{!}{%
\begin{tabular}{lccc}
\toprule
Model & bounding & coverage / class-1 coverage & width \\
\midrule
Frozen-Temporal Head Control & clipping & \textbf{0.8996 / 0.2104} & 0.2685 \\
Frozen-Temporal Head Control & logit & 0.0000 / 0.0000 & 0.0165 \\
Hybrid GRU-D & clipping & \textbf{0.8998 / 0.2227} & 0.2580 \\
Hybrid GRU-D & logit & 0.0000 / 0.0000 & 0.0153 \\
Temporal-Only GRU-D & clipping & \textbf{0.9003 / 0.2153} & 0.3159 \\
Temporal-Only GRU-D & logit & 0.0000 / 0.0000 & 0.0227 \\
\midrule
Latency path & \multicolumn{3}{c}{median / p95 / p99 / max (ms) or amortized per arrival} \\
\midrule
forward pass, GPU & \multicolumn{3}{c}{6.84 / 8.78 / 9.22 / 10.8} \\
forward pass, CPU & \multicolumn{3}{c}{1.58 / 1.93 / 2.17 / 3.2} \\
end to end: GPU forward + ECI step + set & \multicolumn{3}{c}{6.92 / 8.87 / 9.49 / 13.3} \\
split-conformal set & \multicolumn{3}{c}{0.0003 per arrival} \\
ECI step + set & \multicolumn{3}{c}{0.0009 per arrival} \\
Mondrian-ECI step + set & \multicolumn{3}{c}{0.0011 per arrival} \\
\bottomrule
\end{tabular}%
}
\end{table*}


\subsection{NEWS2 baseline sensitivity}

\begin{figure}[H]
\centering
\includegraphics[width=0.85\textwidth]{figures/fig5_news2_sensitivity.pdf}
\caption{Six-arm NEWS2 sensitivity analysis (\texttt{M2\_Findings\_Dossier.md} \S23 row 3).}
\end{figure}

\begin{table}[t]
\centering
\caption{NEWS2 sensitivity analysis on the chronological test period. NEWS2 is computed from 5 of its 7 parameters (consciousness level and supplemental oxygen are not available); a missing parameter scores 0 points. Scored at triage (Arm A), it falls well below the AUROC range reported in published ED validations (0.71-0.86). Arms B--D change the outcome definition. Arm E excludes the 6,302 stays with a missing triage parameter; Arm F rescores on the last-recorded ED vitals and excludes stays with none. Highest AUROC in bold. Proxy cohort (MIMIC-IV-ED-derived).}
\label{tab:m2-news2-sensitivity}
\resizebox{\columnwidth}{!}{%
\begin{tabular}{lccc}
\toprule
Arm & $n$ (positives) & AUROC & 95\% CI \\
\midrule
A: triage, composite outcome & 84,241 (6815) & 0.5507 & [0.5423, 0.5583] \\
B: triage, in-hospital death & 84,241 (1070) & 0.5471 & [0.5265, 0.5659] \\
C: triage, ICU transfer & 84,241 (6500) & 0.5456 & [0.5369, 0.5533] \\
D: triage, death within 24 h of ED departure & 84,241 (143) & 0.4419 & [0.3865, 0.4893] \\
E: triage, complete cases, composite & 77,939 (4394) & 0.6715 & [0.6623, 0.6800] \\
F: last-recorded ED vitals, composite & 83,691 (6495) & \textbf{0.7187} & [0.7114, 0.7261] \\
\bottomrule
\end{tabular}%
}
\end{table}


\subsection{Zero-ED-observation-rate confound check}
Restricting every model's AUROC to test-period stays with $\geq$1 real ED-vitalsign observation (84{,}224 of 84{,}241 stays, 99.98\%) reproduces the full-population numbers in Table~\ref{tab:m2-v2-headline} to within 0.0001--0.0007 AUROC for all 6 models -- no evidence the train/test recording-practice shift (\S\ref{sec:data}) confounds any discrimination result (\texttt{data/m2\_v2\_recording\_confound\_check.json}).

% ======================================================================
\section{Discussion}
% ======================================================================
\begin{prosetodo}{Deferred}
Draws from \texttt{M2\_Findings\_Dossier.md} \S1's framing note (Grinsztajn et al., NeurIPS 2022, on tabular-data skepticism), \S24.4, and \S23 (all four external Q1-readiness blockers, dispositioned).
\end{prosetodo}

% ======================================================================
\section{Limitations}
% ======================================================================

\begin{enumerate}[leftmargin=*]
\item \textbf{The horizon ablation is partial.} Composite-outcome tuned hyperparameters were used at every horizon (not re-tuned), only 3 of 6 models were run (so the static-contribution decomposition is not repeated per horizon), the 1h arm is underpowered and uninterpretable (66 test positives), and no horizon-specific conformal analysis was run (Sec.~\ref{sec:horizon}).
\item \textbf{Secondary (random-split) robustness arm not yet run.} Subject-id partition persisted; execution deliberately deferred (\texttt{M2\_Findings\_Dossier.md} \S24.5) pending a parametrized-caching refactor.
\item \textbf{mTAN's search space is narrower than XGBoost's and GRU-D's.} The matched-rigor re-tune (Deviation 15, resolved) used 3 seeds, 4 stages and 13 candidates with the same noise and calibration guards; the class-imbalance stage does not apply to its shared loss, and the search covered capacity, reference points, heads, learning rate and weight decay only.
\item \textbf{ECI is not needed against periodic recalibration.} In a pre-registered comparison (\S24.3c) periodic recalibration of ordinary split-conformal restores marginal validity with tighter intervals than ECI on all three models; the module's roadmap-stated primary contribution is therefore not supported against the cheap remedy, and the contribution is an empirical characterization with pre-registered negative results. Label lag, the axis most likely to separate the methods, is untested.
\item \textbf{DeLong's test omitted;} the patient-cluster bootstrap is the sole inferential procedure (no validated cluster-adjusted implementation exists).
\item \textbf{The online label-conditional construction (Mondrian-ECI sets) is not a supported contribution.} Against pre-registered comparators, a static Mondrian method with periodic recalibration restores the same coverage with smaller sets (\S24.3b, Part~4 of the pre-registration); the finding that survives is that feature-space conditioning (score-binned groups) fails on the minority class (class-1 coverage 0.09--0.29).
\item \textbf{Label lag is tested only as a uniform delay.} Real resolution times differ by patient and outcome (an ICU transfer resolves quickly, an in-hospital death can take weeks), lag is measured in arrivals with an approximate hour conversion, and label modification by intervention (alerts changing the outcomes later observed) is not modelled (\S\ref{sec:label-lag}).
\item \textbf{NEWS2 baseline is partial} (4 of 7 real clinical-score parameters; no SpO$_2$ scale, supplemental-oxygen term, or consciousness component) and is scored at triage only, which Arm~F's sensitivity analysis shows accounts for the residual gap to published ED-validation ranges.
\end{enumerate}

% ======================================================================
\section{Reproducibility}
% ======================================================================

\texttt{python experiments/run\_reproducible\_m2.py --config configs/paper\_track\_m2.yaml} reproduces every v2 result in this report end to end, in dependency order (cohort rebuild $\to$ tuning $\to$ 5-seed confirmatory runs $\to$ conformal $\to$ inferential statistics $\to$ publication artifacts). Every trained checkpoint, per-seed result JSON, per-encounter prediction, and tuning trace is saved under \texttt{data/}.

% ======================================================================
\section{T=2 (Superseded) Pipeline -- Historical Reference}
\label{sec:t2-legacy}
% ======================================================================

The fixed-T=2 (triage + latest-ED), random-split pipeline that preceded Gate v2 (\texttt{M2\_Findings\_Dossier.md} \S\S1--22) is retained, independently reproducible from \texttt{data/archive\_pre\_irregular/}, and never silently mixed with the v2 results above. \texttt{M2\_Outcome\_Definition\_Gate.md}'s v1 was never signed off; v2 supersedes it in full.

\begin{prosetodo}{Deferred}
If retained in the final manuscript at all (likely as a single retrospective paragraph on methodological evolution, not a parallel result set), drawn from \texttt{M2\_Findings\_Dossier.md} \S\S1--22.
\end{prosetodo}

% ======================================================================
\begin{thebibliography}{9}

\bibitem[Che et al., 2018]{che2018recurrent}
Che, Z., Purushotham, S., Cho, K., Sontag, D., \& Liu, Y. (2018).
Recurrent Neural Networks for Multivariate Time Series with Missing Values.
\emph{Scientific Reports}, 8, 6085.

\bibitem[Shukla \& Marlin, 2021]{shukla2021multi}
Shukla, S.~N., \& Marlin, B. (2021).
Multi-Time Attention Networks for Irregularly Sampled Time Series.
\emph{International Conference on Learning Representations (ICLR)}.

\bibitem[Wu et al., 2025]{wu2025error}
Wu, J., Hu, D., Bao, Y., Xia, S., \& Zou, C. (2025).
Error-quantified Conformal Inference for Time Series.
\emph{International Conference on Learning Representations (ICLR)}.

\bibitem[Sadinle et al., 2019]{sadinle2019least}
Sadinle, M., Lei, J., \& Wasserman, L. (2019).
Least Ambiguous Set-Valued Classifiers with Bounded Error Levels.
\emph{Journal of the American Statistical Association}, 114(525), 223--234.

\bibitem[Quan et al., 2005]{quan2005coding}
Quan, H., Sundararajan, V., Halfon, P., et al. (2005).
Coding Algorithms for Defining Comorbidities in ICD-9-CM and ICD-10 Administrative Data.
\emph{Medical Care}, 43(11), 1130--1139.

\bibitem[Chen and Guestrin, 2016]{chen2016xgboost}
Chen, T., \& Guestrin, C. (2016).
XGBoost: A Scalable Tree Boosting System.
In \emph{Proceedings of the 22nd ACM SIGKDD International Conference on Knowledge Discovery and Data Mining (KDD'16)}, 785--794.

\bibitem[Johnson et al., 2023]{johnson2023mimic}
Johnson, A.~E.~W., Bulgarelli, L., Shen, L., et al. (2023).
MIMIC-IV, a Freely Accessible Electronic Health Record Dataset.
\emph{Scientific Data}, 10, 1.

\bibitem[MAPIE Contributors]{mapie}
MAPIE Contributors. \emph{MAPIE: Model Agnostic Prediction Interval Estimator} [Software]. \url{https://github.com/scikit-learn-contrib/MAPIE}.

\bibitem[Grinsztajn et al., 2022]{grinsztajn2022tree}
Grinsztajn, L., Oyallon, E., \& Varoquaux, G. (2022).
Why do tree-based models still outperform deep learning on tabular data?
\emph{Advances in Neural Information Processing Systems (NeurIPS), Datasets and Benchmarks Track}.

\end{thebibliography}

\end{document}
